\documentclass[11pt]{article}
\usepackage[margin=1in]{geometry}
\usepackage{amsmath, amssymb}
\usepackage{graphicx}
\usepackage{hyperref}
\usepackage{listings}
\usepackage{xcolor}

\definecolor{codegreen}{rgb}{0,0.6,0}
\definecolor{codegray}{rgb}{0.5,0.5,0.5}
\definecolor{codepurple}{rgb}{0.58,0,0.82}
\definecolor{backcolour}{rgb}{0.95,0.95,0.92}

\lstdefinestyle{mystyle}{
    backgroundcolor=\color{backcolour},   
    commentstyle=\color{codegreen},
    keywordstyle=\color{magenta},
    numberstyle=\tiny\color{codegray},
    stringstyle=\color{codepurple},
    basicstyle=\ttfamily\footnotesize,
    breakatwhitespace=false,         
    breaklines=true,                 
    captionpos=b,                    
    keepspaces=true,                 
    numbers=left,                    
    numbersep=5pt,                  
    showspaces=false,                
    showstringspaces=false,
    showtabs=false,                  
    tabsize=2
}
\lstset{style=mystyle}

\title{Hardware-Accelerated Reinforcement Learning Simulator for Octopus CXL Memory Pooling}
\author{Pulak Mehrotra}
\date{\today}

\begin{document}

\maketitle

\begin{abstract}
The current trace-driven Octopus simulator relies on sequential, CPU-bound execution, limiting the sample efficiency of Reinforcement Learning (RL) agents such as Soft Actor-Critic (SAC)[cite: 324, 348]. To effectively train policies capable of anticipating diurnal cloud workloads [cite: 128] and transitioning to online Locally Constrained Policy Optimization (LCPO)[cite: 339], the environment must be highly parallelized. This proposal outlines the transition to a JAX-based, fully tensorized simulation architecture deployed across a 3x TITAN RTX GPU cluster, enabling millions of simulated transitions per second.
\end{abstract}

\section{Phase 1: Stateless Functional Representation (JAX PyTrees)}
Standard object-oriented environment states inherently bottleneck hardware acceleration. In JAX, the simulation state must be immutable and represented as a PyTree. For the target 96-host Acadia pod [cite: 538] connected to $m=120$ Multi-Ported CXL Devices (MPDs)[cite: 538], the global state is defined without batch dimensions, allowing JAX to handle vectorization natively.

\begin{lstlisting}[language=Python, caption=JAX PyTree State Definition]
import jax.numpy as jnp
from flax import struct

@struct.dataclass
class OctopusState:
    mpd_loads: jnp.ndarray       # Shape: (120,) - Load per MPD
    active_vms: jnp.ndarray      # Shape: (Max_VMs, 3) - [host_id, mpd_id, size]
    current_tick: jnp.int32      # Scalar - Time step in the 14-day trace
\end{lstlisting}

\section{Phase 2: Tensorized Environment Dynamics}
The core memory allocation logic must be rewritten as pure linear algebra. When a VM arrives at host $S_i$, the RL policy maps the system state to an allocation vector $a \in \mathbb{R}_{\ge0}^{m}$[cite: 31]. The graph sparsity dictates that $a_j = 0$ if the adjacency matrix $M_{ij} = 0$[cite: 30, 33].

In the JAX environment step, this hard constraint is enforced by applying a logarithmic mask prior to the softmax normalization:

\begin{equation}
    \tilde{w}_j = 
    \begin{cases} 
      w_j & \text{if } M_{ij} = 1 \\
      -\infty & \text{if } M_{ij} = 0 
    \end{cases}
\end{equation}

\begin{equation}
    a_j = x \cdot \frac{\exp(\tilde{w}_j)}{\sum_{k=1}^m \exp(\tilde{w}_k)}
\end{equation}

This operation executes via element-wise matrix multiplication, eliminating iterative masking loops and allowing constant-time execution on the GPU.

\section{Phase 3: Zero-Copy Trace Execution}
The production VM traces from Azure datacenters (e.g., AMS20, LON23) span approximately 14 days at a 5-minute resolution[cite: 107, 127]. To prevent PCIe bandwidth from throttling the simulation, the entire trace dataset is pre-compiled into a static three-dimensional tensor: $\mathcal{T} \in \mathbb{R}^{T \times n \times F}$, where $T = 4032$ (ticks), $n = 96$ (hosts), and $F$ represents features like VM size and lifetime[cite: 279]. 

Prior to training, $\mathcal{T}$ is pinned directly to the TITAN RTX VRAM. During the environment \texttt{step()}, the \texttt{jax.lax.dynamic\_slice} primitive fetches the incoming demand for \texttt{state.current\_tick} entirely within silicon, achieving zero-copy data loading.

\section{Phase 4: Multi-GPU Scaling Strategy (3x TITAN RTX)}
To exploit the 72GB of aggregate VRAM across the 3x TITAN RTX cluster, the single-pod \texttt{step()} function is transformed using JAX's composable parallelization primitives.

Let $B_{gpu}$ denote the maximum number of parallel pods that fit within the 24GB VRAM of a single TITAN RTX (estimated $B_{gpu} \approx 4096$). 

\begin{lstlisting}[language=Python, caption=Hierarchical Parallelization]
import jax

# 1. Intra-GPU Vectorization: Simulates B_gpu pods concurrently
vmapped_step = jax.vmap(step_fn, in_axes=(0, 0))

# 2. Inter-GPU Parallelization: Distributes across 3 GPUs
# Total parallel simulations = 3 * B_gpu
parallel_env_step = jax.pmap(vmapped_step, in_axes=(0, 0))
\end{lstlisting}



This hierarchical scaling ensures that the SAC agent experiences a massively diverse set of diurnal permutations [cite: 128, 131] and randomized link failures [cite: 519] simultaneously, stabilizing the policy gradient and closing the gap against the optimal offline baseline[cite: 524].

\end{document}